\documentclass{beamer}
\usetheme{default}
\useoutertheme{miniframes}
\useinnertheme{circles}
\usepackage{listings}
\usepackage{xcolor}
\usepackage{bookmark}
\usepackage{url}

% =============================================
% COLOR SCHEME: Professional Navy
% =============================================
\definecolor{NavyDark}{RGB}{20, 40, 80}        % Deep navy
\definecolor{NavyMid}{RGB}{40, 65, 115}        % Medium navy
\definecolor{NavyLight}{RGB}{70, 100, 150}     % Lighter navy accent

% Set the main structure color
\setbeamercolor{structure}{fg=NavyDark}

% Header (top navigation bar)
\setbeamercolor{section in head/foot}{fg=white, bg=NavyDark}
\setbeamercolor{subsection in head/foot}{fg=white, bg=NavyMid}

% Footer colors (subtle gradient)
\setbeamercolor{author in head/foot}{fg=white, bg=NavyLight}
\setbeamercolor{title in head/foot}{fg=white, bg=NavyMid}
\setbeamercolor{date in head/foot}{fg=white, bg=NavyDark}

% Title page
\setbeamercolor{title}{fg=NavyDark}
\setbeamercolor{subtitle}{fg=NavyMid}
\setbeamercolor{author}{fg=NavyDark!80}
\setbeamercolor{date}{fg=NavyDark!60}

% Frame titles
\setbeamercolor{frametitle}{fg=white, bg=NavyDark}
\setbeamercolor{framesubtitle}{fg=NavyMid}

% Itemize/enumerate colors
\setbeamercolor{item}{fg=NavyDark}
\setbeamercolor{subitem}{fg=NavyMid}

% Navigation dots
\setbeamercolor{mini frame}{fg=white, bg=white!30}
\setbeamercolor{mini frame inactive}{fg=white!50}

% Custom footline: Author | Short Title | Date + Frame Number
\setbeamertemplate{footline}{
  \leavevmode%
  \hbox{%
  \begin{beamercolorbox}[wd=.20\paperwidth,ht=2.5ex,dp=1ex,left,leftskip=1em]{author in head/foot}%
    \usebeamerfont{author in head/foot}\insertshortauthor
  \end{beamercolorbox}%
  \begin{beamercolorbox}[wd=.55\paperwidth,ht=2.5ex,dp=1ex,center]{title in head/foot}%
    \usebeamerfont{title in head/foot}\insertshorttitle
  \end{beamercolorbox}%
  \begin{beamercolorbox}[wd=.25\paperwidth,ht=2.5ex,dp=1ex,right,rightskip=1em]{date in head/foot}%
    \usebeamerfont{date in head/foot}\insertshortdate{}\hspace*{2em}
    \insertframenumber{} / \inserttotalframenumber
  \end{beamercolorbox}}%
  \vskip0pt%
}

% Code listing style
\lstset{
    basicstyle=\ttfamily\footnotesize,
    breaklines=true,
    frame=single,
    backgroundcolor=\color{gray!10},
}

\title[Systems Thinking Workshop Series 1: Experiment Tracking]{Systems Thinking for Better Science: \\ \textbf{Experiment Tracking Workshop}}
\subtitle{\textbf{How to make your experiments reproducible and shareable}}
\author[Sneha \& Adish]{\textbf{Sneha \& Adish}}
\date[Dec 18, 2025]{\today}

\begin{document}

\maketitle

\begin{frame}{Workshop Series: Systems Thinking for Better Science}
    \textbf{Four Sessions Building on Each Other:}
    
    \vspace{0.5em}
    \begin{enumerate}
        \item \textbf{Session 1 - Experiment Tracking (Today)} \\
        {\small Git workflows, reproducibility, why it matters}
        
        \item \textbf{Session 2 - Documentation Best Practices (Jan 2026)} \\
        {\small Guides vs references, auto-generation, publishing}
        
        \item \textbf{Session 3 - Advanced Experiment Tracking (Jan 2026)} \\
        {\small Branching, migration, complex workflows, secrets management}
        
        \item \textbf{Session 4 - Linear Integration (Feb 2026)} \\
        {\small Connecting experiments to team planning, linear integration}
    \end{enumerate}
    
    \vspace{0.5em}
    \textit{Each session will build on the previous. The intention is to go slow and level up progressively.}
\end{frame}


% ============================================
% SECTION 1: HOOK
% ============================================
\section{Hook}



\begin{frame}{Quick Question}
\Large
Raise your hand if you've experienced any of the following:
\vspace{1em}

\begin{itemize}
    \item You've lost work or can't reproduce previous results
    \pause
    \item A colleague asks to reproduce your analysis - you have an idea of what you did but can't remember exactly, or you're not sure if you did it correctly.
    \pause
    \item You broke something and wished for an undo button.
\end{itemize}
\end{frame}

\begin{frame}{These Aren't Just Frustrations}
\Large
They're symptoms of \textbf{missing systems}. Systems on how to work with data, run experiments, and save and share results.

\vspace{2em}

\normalsize
The good news? We can fix this.

\vspace{2em}

\normalsize
Today, we'll learn how to use git to track experiments, and how to make your experiments reproducible and shareable.
\end{frame}

\begin{frame}{Today's Journey}
\begin{enumerate}
    \item Watch a live experiment workflow
    \item Extract the principles that make it work
    \item Get tools to implement it yourself
    \item See the path to team-wide adoption (March 2026)
\end{enumerate}
\end{frame}

% ============================================
% SECTION 2: DEMO SETUP
% ============================================
\section{Demo}

\begin{frame}[fragile]{Let's Start an Experiment}
    \Large
    Live demo time!
    
    \vspace{1em}
    \normalsize
    We'll build an experiment from scratch using real data from our dataio platform.
    
    \vspace{1.5em}
    \textbf{Follow along later:}
    \url{https://github.com/dsih-artpark/example.git}
    \end{frame}

\begin{frame}[fragile]{Git Commands Reference}
    \begin{description}
        \item[\texttt{git status}] Check current state - see what changed
        
        \item[\texttt{git add <file>}] Stage file for commit \\
        \textcolor{gray}{\textit{Example: \texttt{git add eda.py}}}
        
        \item[\texttt{git commit -m "..."}] Save checkpoint with message \\
        \textcolor{gray}{\textit{Example: \texttt{git commit -m "Add EDA"}}}
        
        \item[\texttt{git log --oneline}] View commit history (your checkpoints)
        
        \item[\texttt{git checkout <hash>}] Go back to a previous commit \\
        \textcolor{gray}{\textit{Example: \texttt{git checkout abc123}}}
        
        \item[\texttt{git diff}] See what changed between versions
    \end{description}
    
    \vspace{1em}
    \centering
    \textit{Keep this visible during demo - you'll see these in action}
    \end{frame}

\begin{frame}[fragile]{The Structure}
\begin{lstlisting}
README.md
.gitignore
.env
.python-version
pyproject.toml
uv.lock
config.yaml
data/
outputs/
src/example/experiments/
    └── 2025-12-18-eda/
        ├── README.md
        ├── __init__.py
        └── eda.py
\end{lstlisting}

\normalsize
This simple structure can help you keep your experiments organised and reproducible.
\end{frame}

\begin{frame}{Demo Plan: What You'll See}
    \textbf{I'll walk through a real workflow:}
    
    \vspace{0.5em}
    \begin{enumerate}
        \item \textbf{Setup:} Repo structure already created
        \item \textbf{First experiment:} Run through various aspects of the experiment
        \item \textbf{Commit:} Save checkpoint with meaningful message
        \item \textbf{Iterate:} Run spatial analysis, calculate coverage
        \item \textbf{Break \& fix:} Make a mistake → recover with git
        \item \textbf{Collaborate:} Show how Adish can reproduce my exact work
        \item \textbf{Linear:} Quick look at organizational layer
    \end{enumerate}
    
    \vspace{0.5em}
    \textit{Watch the workflow, not the code. We'll extract principles after.}
    \end{frame}

% ============================================
% SECTION 3: PRINCIPLES
% ============================================
\section{Lessons Learnt}

\begin{frame}{Lessons Learnt from the Demo}
\Large
Let's work to extract what made that workflow work.
\end{frame}

\begin{frame}[fragile]{Principle 1: Commit Discipline}
\textbf{Commits = Time Machine Checkpoints}

\vspace{1em}
When to commit:
\begin{itemize}
    \item After each meaningful change
    \item Not every keystroke, not once a month
\end{itemize}

\vspace{1em}
Good commit message:
\begin{lstlisting}
Increased window size to 14 days - RMSE dropped to 0.82
\end{lstlisting}

\vspace{0.5em}
You don't need to push every time you commit! So you can make commits locally, roll them back, anbd push them at the end of the day. Committing is a checkpoint, not a push.
\end{frame}

\begin{frame}{Principle 2: Structure Discipline}
\textbf{Reproducible Experiments by Design}

\vspace{1em}
Key elements:
\begin{itemize}
    \item Config-driven (no hardcoded parameters)
    \item Consistent naming (\texttt{experiments/description/})
    \item README documents intent
    \item Separation: data / experiments / results
\end{itemize}

\vspace{1em}
\textit{Structuring your work brings about some clarity of thought, and enables reproducibility and hence collaboration.}
\end{frame}

\begin{frame}[fragile]{Principle 3: Recovery Discipline}
\textbf{Fearless Iteration and Recovery}

\vspace{1em}
You saw:
\begin{lstlisting}
# Made it worse? No problem.
git checkout <previous-commit>
# Back to good state in 3 seconds
\end{lstlisting}

\vspace{1em}
This means:
\begin{itemize}
    \item You can always roll back to a previous checkpoint - so experiment away!
    \item Commit and/or tag important milestones so it's easier to roll back to a previous checkpoint
    \item Don't worry about breaking things
\end{itemize}
\end{frame}

\begin{frame}{Systems Thinking → Better Science}
\Large
These principles matter beyond git:

\vspace{1em}
\normalsize
\begin{itemize}
    \item Clear mental model of your work - you know what you've done, and what you've not done
    \item Codebase is navigatable, for yourself and for others; Less time spent firefighting and reproducing experiments
    \item Better collaboration with tools (LLMs) while structural decisions come from you; HITL stays central
    \item More time for actual thinking, less firefighting
\end{itemize}

\vspace{1em}
\textit{Chat with LLMs about your experiments. But don't blindly hand off to them.}
\end{frame}

\begin{frame}{The Cherry on Top: Linear}
\textbf{Organizational Layer}

\vspace{1em}
\begin{itemize}
    \item Link commits to work tickets
    \item Track progress across team
    \item Integrate with planning cycles
\end{itemize}

\vspace{1em}
\textit{We'll dive deeper in a future workshop.}
\end{frame}

% ============================================
% SECTION 4: NEXT STEPS
% ============================================
\section{Next Steps}

\begin{frame}{What You Get Today}
    \textbf{Resources:}
    \begin{itemize}
        \item \textbf{SoP:} Will be published to Knowledge Base. You can refer to this repo as an example, will be notified when it's ready.
        \item \textbf{Repo template:} A cookiecutter based template will be available in Jan 2026. Until then, you can use this repo as a template.
        \item \textbf{Git cheat sheet:} In this presentation
    \end{itemize}
    
    \vspace{1em}
    \textbf{Support:}
    \begin{itemize}
        \item \textbf{Office hours:} Tue 3-5pm, Thu 11-1pm
        \item Don't struggle alone - come with questions, or to discuss your ideas!
    \end{itemize}

    \vspace{1em}
    \small
    \textit{Adish, Prerna, Tarun,and Aishwarya can attest - I am always happy to sit and chat, answer questions, and help you out.}
    \end{frame}

    \begin{frame}{Your Action Item for Monday}
        \Large
        Start small - one experiment.
        
        \vspace{1em}
        \normalsize
        \textbf{Step 1:} Create ONE experiment repo using the template
        
        \textbf{Step 2:} Run ONE analysis with git commits after each step
        
        \textbf{Step 3:} Bring questions to Tuesday office hours (3-5pm).
        
        \vspace{2em}
        \large
        With git specifically, and structured thinking in general, the key is to avoid overthinking it and start small.
        \end{frame}

\begin{frame}{Your Timeline: Now → March 2026}
\textbf{This Month (Dec 2025):}
\begin{itemize}
    \item Try this workflow for ONE experiment
    \item Come to office hours if stuck
    \item Share what worked/didn't in Slack (\href{https://artpark-in.slack.com/archives/C0638EBD84D}{#technical-support})
\end{itemize}

\vspace{1em}
\textbf{March 2026:}
\begin{itemize}
    \item Official policy for dengue team (then org-wide) - will be notified when it's ready.
    \item Code review gates will be created.
    \item Linear integration required
\end{itemize}

\vspace{1em}
\textit{Early adopters have 3 months of practice. \\
Get ahead of the curve!}
\end{frame}

\begin{frame}{Next Workshops}
\textbf{Session 2:} Documentation Best Practices
\begin{itemize}
    \item Guides vs references
    \item Auto-generate from git history
    \item Tools: Sphinx, MyST, Starlight, GitHub Pages
\end{itemize}

\vspace{1em}
\textbf{Session 3:} Advanced Experiment Tracking
\begin{itemize}
    \item Branching strategies
    \item Migrating existing messy projects
    \item Complex collaboration workflows
    \item Secrets management using 1Password
\end{itemize}

\vspace{1em}
\textbf{Session 4:} Deep dive into Linear integration
\begin{itemize}
    \item Connecting experiments to team planning
    \item Linear integration with pull requests
\end{itemize}
\end{frame}

\begin{frame}{The Alternative: Life Without This System}
    \begin{columns}
    \column{0.48\textwidth}
    \textbf{Without Git \& Structure:}
    \begin{itemize}
        \item Lost experiments
        \item "Which parameters worked?"
        \item Can't reproduce for papers
        \item Fear of breaking things
        \item Firefighting chaos
    \end{itemize}
    
    \column{0.48\textwidth}
    \textbf{With This System:}
    \begin{itemize}
        \item Time machine for work
        \item Clear experiment history
        \item Reproducible science
        \item Fearless iteration
        \item Focus on thinking
    \end{itemize}
    \end{columns}
    
    \vspace{1em}
    \centering
    \textit{This isn't extra work. It actually helps you reduce firefighting and focus on thinking in the long run.}
    \end{frame}

\begin{frame}{Remember}
\Large
This isn't about rules.

\vspace{1em}
\normalsize
It's about \textbf{self-preservation} and \textbf{better data science}.

\end{frame}

\begin{frame}{Questions?}
\begin{center}
\Large
Let's discuss!

\vspace{2em}
\normalsize
Sneha's Office Hours: Tue 3-5pm, Thu 11-1pm \\
Adish's Office Hours: Ad-hoc, by appointment \\
Reach out! We're here to help.
\end{center}
\end{frame}

\end{document}